\chapter{حماية الجهاز العصبي}\label{ch:g8:protecting-the-nervous-system}

أظهر فصلان ما \omterm{def:g5:brain-in-command:system}{الجهاز العصبي}: أسلاك تتعلم وصلاتها وتقرر
--- ويمكن خداعها
(\cref{rem:g8:nervous-communication:vulnerable}). وختام
هذه السنة يحوّل ذلك التشريح إلى حماية. فماذا يفعل نقص
\omterm{prop:g1:healthy-habits:needs}{النوم} والضجيج والشاشات والمواد المؤثرة في العقل، وصلةً
وصلة، بالآلة التي هي أنت --- وكيف تبدو حمايتها عمليًا،
في الرابعة عشرة؟

\section{الصيانة: النوم، وحدود الحواس}

\begin{proposition}[النوم صيانة للمشابك]\label{prop:g8:protecting-the-nervous-system:sleep}
لحفظ الملفات والإصلاحات في
\cref{prop:g5:brain-in-command:protect} قراءة على مستوى
الوصلات: ففي \omterm{prop:g1:healthy-habits:needs}{النوم} ترسّخ \omterm{def:g8:nervous-communication:synapse}{مشابك} النهار المقواة ---
والتمرين الذي يجعل الأمر دائمًا
(\cref{ex:g8:nervous-communication:juggler}) يحدث ليلًا
في أكثره --- ويزيل تدبير \omterm{def:g5:brain-in-command:system}{الدماغ} المنزلي بلى النهار
الكيميائي. والليالي القصيرة تترك تعلّم الأمس نصف محفوظ
ووزن الغد متثاقلًا: \omterm{def:g5:brain-in-command:system}{والدماغ} \omterm{def:g3:stages-of-life:stages}{المراهق}، وهو في وسط إعادة
التوصيل (\cref{rem:g8:puberty:moods}) وقد انزاحت ساعته
متأخرة، أحوج ما يكون إلى ساعاته حين تبدأ المدرسة أبكر
ما تبدأ.
\end{proposition}
\admitted

\begin{example}[مفارقة المراجعة]\label{ex:g8:protecting-the-nervous-system:revision}
تلميذان قبل اختبار: أحدهما يراجع إلى منتصف الليل ثم إلى
الثانية؛ والآخر يراجع إلى العاشرة وينام. والثاني يفوز في
العادة --- فالمادة عبرت وصلاته \emph{ورسّخت}؛ أما ساعات
الأول الزائدة فكلفته الحفظ نفسه الذي يثبّت المراجعة.
و``نم عليها'' هندسة \omterm{def:g8:nervous-communication:synapse}{مشابك}.
\end{example}

\begin{proposition}[الصخب جرعة]\label{prop:g8:protecting-the-nervous-system:noise}
\omterm{def:g5:first-look-at-cells:cell}{خلايا} \omterm{def:g1:the-five-senses:sense}{الأذن} الحسية --- المتلقيات الشعرية الدقة التي تحول
الصوت إلى رسائل عصبية --- لا تنمو من جديد
(\cref{prop:g1:the-five-senses:fragile} قالها برفق؛ وهذه
هي الآلية). والصوت الصاخب يقتلها بالجرعة: شدة في مدة.
فالحفلات وسماعات \omterm{def:g1:the-five-senses:sense}{الأذن} على أعلى شدة تنفق، ساعة بعد
ساعة، رأس مال من المتلقيات لا يرده \omterm{prop:g1:healthy-habits:needs}{نوم} --- والصفير بعد
ليلة صاخبة (\cref{ex:g1:the-five-senses:loud}) هو تحذير
الجرعة الزائدة. والقاعدة ميزانية: أخفض، وأقصر، وبفواصل
--- وسدادات \omterm{def:g1:the-five-senses:sense}{أذن} حيث تحكم المكبرات.
\end{proposition}
\admitted

\section{المواد التي تخدع الوصلات}

\begin{proposition}[المواد المؤثرة في العقل تعمل عند المشابك]\label{prop:g8:protecting-the-nervous-system:substances}
كل مادة تغير المزاج أو الإدراك أو اليقظة تعمل عند فجوات
\cref{def:g8:nervous-communication:synapse} --- فتحاكي
رسولًا، أو تسده، أو تفرض إطلاقه:
\begin{itemize}
  \item \emph{الكحول} يثبط وصلات الجهاز كله: فتبطؤ
        الحلقة من طرفها إلى طرفها --- الحكم، \omterm{def:g5:brain-in-command:reaction}{وزمن رد
        الفعل} (فتطول أمتار
        \cref{ex:g5:brain-in-command:driver})، والتوازن،
        والذاكرة؛ ووصلات \omterm{def:g5:brain-in-command:system}{الدماغ} \omterm{def:g3:stages-of-life:stages}{المراهق} التي لا تزال
        تبنى تأخذ التثبيط أعمق ما يكون؛
  \item \emph{نيكوتين التبغ} يحاكي رسولًا طبيعيًا، وإذا
        تكرر أعاد توصيل وصلات المكافأة نحو
        \emph{\omterm{prop:g5:healthy-choices:substances}{الإدمان}} (فخ
        \cref{prop:g5:healthy-choices:substances}، ومعه
        الآن آليته: \omterm{def:g8:nervous-communication:synapse}{المشابك} تتكيف لتتوقع الجرعة ---
        وتحتج من غيرها)؛
  \item \emph{القنّب} يحاكي أسرة أخرى من الرسل، فيعكّر
        بالضبط وصلات الذاكرة والدافع والتوقيت --- وأكثر
        ما يكلف، مرة أخرى، دماغًا قيد الإنشاء؛
  \item وكلما اشتدت المادة كان فرضها على الوصلة أفظّ ---
        وكانت إعادة توصيل المكافأة، وهي محرك \omterm{prop:g5:healthy-choices:substances}{الإدمان}،
        أسرع.
\end{itemize}
\end{proposition}
\admitted

\begin{remark}[الإدمان، آليًا]\label{rem:g8:protecting-the-nervous-system:addiction}
لفخ \Cref{rem:g5:healthy-choices:never} مخطط توصيل الآن:
فوصلات المكافأة، إذا جرّعت مرارًا، \emph{تتكيف} ---
فتقلّم المستقبلات، وترفع العتبات --- حتى تلزم المادة
ليشعر المرء بأنه عادي، ويصير غيابها احتجاجًا من الجهاز.
وذلك \omterm{def:g4:adapted-to-their-world:adaptation}{التكيف} هو تباين سهولة الدخول وصعوبة الخروج مصنوعًا
من \omterm{def:g8:nervous-communication:synapse}{مشابك}: فالبدء قرار واحد؛ أما التوقف فضد وصلات أعيد
توصيلها. \omterm{def:g5:brain-in-command:system}{والدماغ} \omterm{def:g3:stages-of-life:stages}{المراهق}، وهو ما يزال قابلًا للضبط، أسرع
إعادة توصيل --- ولهذا تجد كل إحصائية المبتدئين باكرًا
أعمق وقوعًا في الفخ، ولهذا يبقى ``ألا تبدأ قط'' هو
الحركة الرخيصة.
\end{remark}

\begin{example}[حلقة الشاشة]\label{ex:g8:protecting-the-nervous-system:screens}
لا جزيء، والوصلات هي هي: فالخلاصات والألعاب مضبوطة ---
عن عمد --- لتدغدغ أسلاك المكافأة بجوائز صغيرة لا تتوقع،
وهو الجدول الذي يدرّب الوصلات أشد تدريب. والنتيجة ابن عم
لطيف للفخ: تفقّد يقاوم التوقف، وأمسيات تتبخر، \omterm{prop:g1:healthy-habits:needs}{ونوم} يدفع
متأخرًا (ساعات
\cref{prop:g5:healthy-choices:screens}، مضافًا إليها
كلفة
\cref{prop:g8:protecting-the-nervous-system:sleep}).
والدفاع هو الدفاع دائمًا: قرر الساعات سلفًا، وأبق الساعة
الأخيرة بلا شاشة.
\end{example}

\section{التطبيق: أن تدير حمايتك بنفسك}

\begin{method}[دليل مالك الجهاز العصبي]\label{met:g8:protecting-the-nervous-system:manual}
قائمة العناية في \cref{met:g5:healthy-choices:run}، وقد
أعيد إصدارها بآلياتها:
\begin{enumerate}
  \item \emph{نم الساعات} --- فالترسيخ والتدبير المنزلي
        يجريان ليلًا؛ واحرس الساعة الأخيرة من الشاشات؛
  \item \emph{وازن ميزانية الصخب} --- فرأس مال
        المتلقيات لا ينمو من جديد: اخفض الشدة، وأدخل
        الفواصل، وضع سدادات \omterm{def:g1:the-five-senses:sense}{الأذن} حيث يزأر؛
  \item \emph{والخوذ} حيث يكون السقوط سريعًا
        (\cref{ex:g5:brain-in-command:helmet}) ---
        فالأسلاك لا تلتئم \omterm{def:g3:bones-and-muscles:skeleton}{كالعظم}؛
  \item \emph{ولا خادعات للوصلات} \omterm{def:g5:brain-in-command:system}{لدماغ} قيد الإنشاء ---
        و``لا'' يقرر هادئًا سلفًا (جملة الجيب في
        \cref{ex:g5:healthy-choices:answers}، ومعها الآن
        تبريرها الكامل)؛
  \item \emph{واستعمل الآلة} --- التعلم والرياضة
        والموسيقى: فالوصلات تقوى بالمرور، والحلقة
        المدربة هي الأصل الذي تحميه كل عادة أخرى.
\end{enumerate}
\end{method}

\begin{example}[المساعدة موجودة، وهي تنفع]\label{ex:g8:protecting-the-nervous-system:help}
ولمن وقع في الشرك سلفًا --- في عادته أو عادة صديق، مواد
كانت أم شاشات: عناوين \cref{ex:g8:contraception:help}
تخدم هنا أيضًا. فممرضات المدارس والأطباء يعاملون \omterm{prop:g5:healthy-choices:substances}{الإدمان}
بوصفه ما يبينه هذا الفصل --- حالة توصيل، لا حكمًا على
الأخلاق --- والوصلات المعاد توصيلها تستطيع أن تعود
فتتكيف: أصعب من ألا تبدأ قط، وبعيد جدًا عن اليأس. وأسوأ
الخطط خطة ساحة المدرسة: الصمت.
\end{example}

\section{تمارين}

\begin{exercise}[$\star$]\label{exo:g8:protecting-the-nervous-system:1}
ماذا يفعل \omterm{prop:g1:healthy-habits:needs}{النوم} عند الوصلات، ومن أحوج الناس إلى
الساعات؟
\end{exercise}

\begin{exercise}[$\star$]\label{exo:g8:protecting-the-nervous-system:2}
لماذا كان الصوت الصاخب مسألة ميزانية؟ ما الذي ينفق، وما
تحذير الجرعة الزائدة؟
\end{exercise}

\begin{exercise}[$\star$]\label{exo:g8:protecting-the-nervous-system:3}
أين تعمل المواد المؤثرة في العقل جميعًا، وبأي طرق ثلاث؟
\end{exercise}

\begin{exercise}[$\star$]\label{exo:g8:protecting-the-nervous-system:4}
ماذا يفعل الكحول بالحلقة --- سمّ ثلاث محطات.
\end{exercise}

\begin{exercise}[$\star$]\label{exo:g8:protecting-the-nervous-system:5}
اذكر آلية \omterm{prop:g5:healthy-choices:substances}{الإدمان} في جملتين.
\end{exercise}

\begin{exercise}[$\star$]\label{exo:g8:protecting-the-nervous-system:6}
لماذا تنتمي حلقة الشاشة إلى هذا الفصل، وهي بلا جزيء؟
\end{exercise}

\begin{exercise}[$\star\star$]\label{exo:g8:protecting-the-nervous-system:7}
فسّر مفارقة المراجعة بمساعدة
\cref{prop:g8:protecting-the-nervous-system:sleep}.
\end{exercise}

\begin{exercise}[$\star\star$]\label{exo:g8:protecting-the-nervous-system:8}
لماذا تكلف الجرعات نفسها \omterm{def:g5:brain-in-command:system}{دماغ} \omterm{def:g3:stages-of-life:stages}{المراهق} أكثر مما تكلف
\omterm{def:g5:brain-in-command:system}{دماغ} البالغ؟ آليتان
(\cref{rem:g8:puberty:moods}،
\cref{rem:g8:protecting-the-nervous-system:addiction}).
\end{exercise}

\begin{exercise}[$\star\star$]\label{exo:g8:protecting-the-nervous-system:9}
اربط أمتار \cref{ex:g5:brain-in-command:driver} بتثبيط
الكحول: ما الذي يطول، ولماذا كان قانون الشرب والقيادة
قانون زمن رد فعل؟
\end{exercise}

\begin{exercise}[$\star\star$]\label{exo:g8:protecting-the-nervous-system:10}
أعد تبرير ثلاثة بنود من دليل المالك بآلياتها على مستوى
الوصلات.
\end{exercise}

\begin{exercise}[$\star\star$]\label{exo:g8:protecting-the-nervous-system:11}
``حالة توصيل، لا حكمًا على الأخلاق.'' ماذا يغير هذا
التأطير في طلب المساعدة --- وفي عرضها على صديق؟
\end{exercise}

\begin{exercise}[$\star\star\star$]\label{exo:g8:protecting-the-nervous-system:12}
اكتب حجة السنة الختامية: من الفجوات القابلة للضبط في
\cref{def:g8:nervous-communication:synapse}، استنبط أعظم
قوة في \omterm{def:g5:brain-in-command:system}{الدماغ} ومواطن ضعفه المميزة معًا --- التعلم
\omterm{prop:g5:healthy-choices:substances}{والإدمان} آلية واحدة تقرأ مرتين.
\end{exercise}

\section{مسألة: أمسية الوقاية}

\begin{problem}[{مسألة نهاية الأسبوع --- أمسية المدرسة
لأولياء الأمور والتلاميذ، بنص مكتوب
بالآليات}]\label{pb:g8:protecting-the-nervous-system:1}
يعد الصف أمسية الوقاية في المدرسة: أربع محطات ---
\omterm{prop:g1:healthy-habits:needs}{النوم}، والصوت، والمواد، والشاشات --- لكل واحدة عرض
وبطاقة آلية.

\medskip\noindent\textbf{الجزء الأول --- بطاقات
المحطات.}
\begin{enumerate}
  \item اكتب \omterm{prop:g3:balanced-diet:jobs}{بطاقة} \omterm{prop:g1:healthy-habits:needs}{النوم}: العرض درجتا متطوعين في اختبار
        المسطرة (\cref{met:g5:brain-in-command:ruler})،
        مرتاح في مقابل قصير الليلة. تنبأ بالدرجتين، وأعط
        آلية الوصلات.
  \item اكتب \omterm{prop:g3:balanced-diet:jobs}{بطاقة} الصوت: مقياس ديسيبل، وزوج سدادات
        \omterm{def:g1:the-five-senses:sense}{أذن}، وجملة ``لا تستطيع أن تنميها من جديد''. املأ
        الآلية.
  \item اكتب \omterm{prop:g3:balanced-diet:jobs}{بطاقة} المواد: مخطط واحد \omterm{def:g8:nervous-communication:synapse}{لمشبك}، وثلاثة أسهم
        عليها: حاكِ، وسد، وافرض. أسند الكحول والنيكوتين
        وسمًا من \cref{exo:g8:nervous-communication:11}
        إلى أسهمها.
  \item اكتب \omterm{prop:g3:balanced-diet:jobs}{بطاقة} الشاشات: لا جزيء --- فما الذي يشترك
        فيه عرض الهاتف ذي الإشعارات مع محطة المواد؟
\end{enumerate}

\medskip\noindent\textbf{الجزء الثاني --- أسئلة
الحضور.}
\begin{enumerate}[resume]
  \item يقول أحد الأولياء: ``شربنا في سنهم وخرجنا
        بخير.'' أجب بتباين ورشة البناء --- ما الذي
        يختلف في الرابعة عشرة؟
  \item ويقول تلميذ: ``سيجارة واحدة لا تعيد توصيل
        شيء.'' أجب بمساعدة
        \cref{rem:g8:protecting-the-nervous-system:addiction}
        و\cref{rem:g5:healthy-choices:never} --- ما ثمن
        الجرعة الأولى الحقيقي؟
  \item ويقول ولي أمر: ``أوليست الشاشات تلفازًا
        جديدًا؟'' فرّق بين الجدولين: ما الذي يجعل
        الجوائز الصغيرة غير المتوقعة أشد تدريبًا؟
  \item ويسأل تلميذ، بهدوء: ``وإن كان أحدهم لا يستطيع
        التوقف أصلًا؟'' أعط جواب
        \cref{ex:g8:protecting-the-nervous-system:help}،
        بعناوينه.
\end{enumerate}

\medskip\noindent\textbf{الجزء الثالث --- الكلمة
الختامية.}
\begin{enumerate}[resume]
  \item تفتتح الكلمة بالمشعوذ وتختتم \omterm{prop:g5:healthy-choices:substances}{بالإدمان} ---
        ``القابلية للضبط نفسها، مقروءة مرتين''. اكتب
        تلك الحجة في أربع جمل.
  \item ثم تسلّم كل تلميذ جملة الجيب في
        \cref{ex:g5:healthy-choices:answers}. فسّر، مع
        الآلية، لماذا يجب أن تؤلف الجملة \emph{الآن}.
  \item يعرض ملصق الأمسية بنود دليل المالك الخمسة.
        رتبها لأسبوع ابن الرابعة عشرة، الأشد حماية
        أولًا، ودافع عن اختيارك الأول.
  \item اختم الأمسية بجملة واحدة: ما الذي يحمى، ولماذا
        كانت الرابعة عشرة وسط العقد الحاسم؟
\end{enumerate}
\end{problem}
